\subsection{A Short History of NLP}
\label{sec:nlp-history}
\slidesrc{Terminology introduction: 4--8}

NLP has gone through five main eras. Each changed \emph{where the knowledge about
language comes from}: first from people writing rules, then from statistics over
corpora, then from neural networks learning their own representations.

\begin{center}
\begin{tikzpicture}[>=Stealth, x=1.6cm]
  \draw[->, thick] (0,0) -- (9.3,0);
  \foreach \x/\lab in {0/1950, 2.6/1990, 5.6/2013, 6.8/2017, 8.2/2022} {
    \draw (\x,0.12) -- (\x,-0.12) node[below, font=\scriptsize] {\lab};
  }
  \fill[nbblue!25]   (0,0.25)   rectangle (2.6,0.85);
  \fill[nbgreen!25]  (2.6,0.25) rectangle (5.6,0.85);
  \fill[nborange!30] (5.6,0.25) rectangle (6.8,0.85);
  \fill[nbpurple!25] (6.8,0.25) rectangle (8.2,0.85);
  \fill[nbred!25]    (8.2,0.25) rectangle (9.2,0.85);
  \node[font=\scriptsize\bfseries] at (1.3,0.55) {Symbolic};
  \node[font=\scriptsize\bfseries] at (4.1,0.55) {Statistical};
  \node[font=\scriptsize\bfseries] at (6.2,0.55) {Neural};
  \node[font=\scriptsize\bfseries, align=center] at (7.5,0.55) {Transf.};
  \node[font=\scriptsize\bfseries] at (8.7,0.55) {LLMs};
  \node[font=\scriptsize, align=center, above] at (1.3,0.9) {rules written\\by hand};
  \node[font=\scriptsize, align=center, above] at (4.1,0.9) {probabilities\\from corpora};
  \node[font=\scriptsize, align=center, above] at (6.2,0.9) {learnt\\features};
  \node[font=\scriptsize, align=center, above] at (7.5,0.9) {attention,\\pre-training};
  \node[font=\scriptsize, align=center, above] at (8.7,0.9) {generative\\chat};
\end{tikzpicture}
\captionof{figure}{The eras of NLP. Boundaries are approximate and the eras overlap: the
neural era is usually dated from word embeddings (around 2013), the transformer era
from the 2017 transformer paper, the LLM era from the release of ChatGPT (2022).}
\end{center}

\subsubsection{Symbolic NLP (1950s -- early 1990s)}
\slidesrc{Terminology introduction: 4}

Programs based on \textbf{hand-written rules} and knowledge bases. Classic systems:
\begin{itemize}[nosep]
  \item \textbf{ELIZA} (1966, Joseph Weizenbaum, MIT): a chatbot imitating a
        psychotherapist, using simple pattern matching and substitution
        (``I am sad'' $\to$ ``Why are you sad?''). It understood nothing, yet many users
        attributed real understanding to it.
  \item \textbf{SHRDLU} (1968--1970, Terry Winograd, MIT): a program that understood
        commands in English about a small ``blocks world'' (``put the red pyramid on the
        blue block'') and could answer questions about it.
  \item \textbf{PARRY} (1972, Kenneth Colby, Stanford): a chatbot simulating a person with
        paranoid schizophrenia; famously, it was made to ``talk'' with ELIZA.
  \item \textbf{Ontologies} (e.g.\ MARGIE, 1973, Roger Schank): representing the
        meaning of sentences with structured conceptual knowledge.
\end{itemize}
Links: \url{https://en.wikipedia.org/wiki/ELIZA}, \url{https://en.wikipedia.org/wiki/SHRDLU},
\url{https://en.wikipedia.org/wiki/PARRY}, \url{https://www.ijcai.org/Proceedings/73}.

\subsubsection{Statistical NLP (1990s -- 2010s)}
\slidesrc{Terminology introduction: 5}

\begin{itemize}[nosep]
  \item \textbf{Large corpora} became available, also multilingual (e.g.\ parliament
        proceedings translated into many languages).
  \item Statistical models make \textbf{soft, probabilistic decisions} by attaching
        real-valued weights to each input feature, instead of applying hard rules.
  \item Example: \textbf{part-of-speech tagging with Hidden Markov Models} (covered
        later in the course).
\end{itemize}

\subsubsection{Neural NLP (2010s)}
\slidesrc{Terminology introduction: 6}

\begin{itemize}[nosep]
  \item \textbf{End-to-end learning} of a higher-level task (e.g.\ question answering),
        instead of relying on a pipeline of separate intermediate tasks (tokenise, tag,
        parse, \ldots).
  \item No need to \textbf{engineer features} by hand: the network learns its own
        representation of the language.
  \item A new paradigm, distinct from statistical NLP (example: \textbf{neural machine
        translation}).
  \item Example: \textbf{word embeddings}, vectors representing words so that words
        with similar meaning are close in space.
\end{itemize}

\subsubsection{Transformers (``it seems a long time ago\ldots'')}
\slidesrc{Terminology introduction: 7}

\textbf{BERT} (Bidirectional Encoder Representations from Transformers, Google) is a
model pre-trained on huge amounts of text and then fine-tuned for specific tasks. It is
built on the \textbf{transformer} architecture, introduced in the paper \emph{Attention
Is All You Need} (Vaswani et al., NeurIPS 2017).

\begin{nbwarning}[dates and the link on the slide]
The slide says ``BERT, Google (2019)'' and links to a NeurIPS 2017 paper. That link is the
\emph{transformer} paper (\emph{Attention Is All You Need}), not the BERT paper. BERT
(Devlin et al.) was released in October 2018 and published at NAACL in 2019, so both
2018 and 2019 are found in the literature.
\end{nbwarning}

\subsubsection{Generative AI and LLMs (present)}
\slidesrc{Terminology introduction: 8}

\begin{itemize}[nosep]
  \item \textbf{GPT}: Generative Pre-trained Transformer. \textbf{ChatGPT} was released in
        November 2022.
  \item Other large language models and tools: Llama, Gemini, \ldots; and Ollama, a tool
        to run open models locally.
\end{itemize}
